\subsection*{章末確認問題}

\addcontentsline{toc}{subsection}{8 章末確認問題}

\begin{itemize}[leftmargin=2em]

\item 専門職実務とは何か。単なるプログラミング能力とどう違うかを説明せよ。

\item 認定、資格証明、資格、免許の違いを、対象と強制力に注目して説明せよ。

\item 倫理規程違反が作為だけでなく不作為でも起こる理由を説明せよ。

\item 専門職団体がソフトウェア技術者にとって重要な理由を三つ挙げよ。

\item 標準に従うことが、品質面と法的面で役立つ理由を説明せよ。

\item 秘密保持契約と知的財産契約で確認すべき点を述べよ。

\item 商標、特許、著作権、営業秘密の違いを一つずつ例を挙げて説明せよ。

\item トレードオフ分析で、評価基準と重みづけを明示する必要がある理由を説明せよ。

\item チーム人数が増えるとコミュニケーションが難しくなる理由を、経路数の観点から説明せよ。

\item 不確実性と曖昧さの違いを説明し、それぞれへの対処を述べよ。

\item 公平性・多様性・包摂性が、分散ソフトウェア開発で重要になる理由を説明せよ。

\item 「読む・理解する・要約する」「書く」「発表する」が、ソフトウェア技術者の仕事にどう関係するかを説明せよ。

\end{itemize}

\begin{center}
\fbox{\parbox{0.94\linewidth}{解答の方向性\par 専門職実務は、知識、技能、態度、倫理、標準、説明責任を含む。認定は教育機関やプログラム、資格証明と資格は個人の能力、免許は法的な業務権限に注目する。標準は最低限の共通土台と説明根拠になる。トレードオフ分析では、目的、評価基準、重み、利害関係者との合意、利益相反の開示が重要である。}}
\end{center}
